\chapter{あとがき}

ここでは、私の主観として何点か述べます。

「鉄屑の騎士と永遠の春」は、短編ではなく長編向けの物語です。
元々短編として制作していましたが、完全な異世界もの（ハイファンタジー）であるため、観客に世界観を理解してもらうまで感情移入の成立は難しいです。
そのため、この物語を採用するとしたら物語プロットを再構成し、数十分の映像作品にする方が良いと考えます。
その場合、過去の制作状況を踏まえると、制作期間は1年ではなく2年かかると思われます。
一応、短編の物語として成立しているので、現在の物語プロットで制作することも可能です。

「約束の果実」は、短編の物語としての完成度は高いと思います。
この物語はローファンタジーであり、異世界ではあるものの控えめなものとなっています。
また、扱うテーマが「家族愛」であり、観客の感情移入が容易であることが最大の理由です。
